% ============================================================
% Shared preamble for every standalone diagram in this folder.
%
% Usage, from a file in this same directory:
%
%   \documentclass[tikz,border=8pt]{standalone}
%   % ============================================================
% Shared preamble for every standalone diagram in this folder.
%
% Usage, from a file in this same directory:
%
%   \documentclass[tikz,border=8pt]{standalone}
%   % ============================================================
% Shared preamble for every standalone diagram in this folder.
%
% Usage, from a file in this same directory:
%
%   \documentclass[tikz,border=8pt]{standalone}
%   % ============================================================
% Shared preamble for every standalone diagram in this folder.
%
% Usage, from a file in this same directory:
%
%   \documentclass[tikz,border=8pt]{standalone}
%   \input{diagramstyle}
%   \begin{document} \begin{tikzpicture} ... \end{tikzpicture} \end{document}
%
% The palette is the one introduced by diagram-orbis.tex. That file
% still carries its own inline copy of these definitions; if you ever
% touch the palette, change it in both places or switch orbis over to
% \input{diagramstyle} as well.
% ============================================================

\usepackage{amsmath,amssymb}
\usepackage{tikz}
\usetikzlibrary{arrows.meta, positioning, calc, fit, backgrounds,
                shapes.geometric, matrix, decorations.pathreplacing}

% ---------- fills -------------------------------------------
% Panel fills are OPAQUE pre-blended colours rather than transparent
% ones, so a line can be given a halo in exactly the panel colour.
% That is what turns an unavoidable crossing into a clean "hop".
\definecolor{panelBlue}{RGB}{248,251,255}
\definecolor{panelGreen}{RGB}{249,253,249}
\definecolor{panelWarm}{RGB}{255,251,247}

\definecolor{semantic}{RGB}{229,224,250}
\definecolor{detail}{RGB}{218,246,228}
\definecolor{latent}{RGB}{255,244,210}
\definecolor{vq}{RGB}{244,232,255}
\definecolor{flow}{RGB}{255,231,222}
\definecolor{attention}{RGB}{224,239,255}
\definecolor{decoder}{RGB}{231,236,255}
\definecolor{neutral}{RGB}{243,244,246}

% ---------- strokes -----------------------------------------
\definecolor{blueDark}{RGB}{25,73,115}
\definecolor{greenDark}{RGB}{35,120,78}
\definecolor{purpleDark}{RGB}{105,70,145}
\definecolor{redDark}{RGB}{145,65,55}
\definecolor{grayDark}{RGB}{78,82,90}
\definecolor{orangeDark}{RGB}{165,95,40}

% ---------- styles ------------------------------------------
\tikzset{
    >=Latex,
    % --- connectors ---
    arrow/.style={->, line width=0.8pt, shorten >=1.5pt, shorten <=1.5pt},
    thinarrow/.style={->, line width=0.6pt, shorten >=1.2pt, shorten <=1.2pt},
    plainline/.style={line width=0.8pt},
    dashedarrow/.style={->, dashed, line width=0.8pt,
                        shorten >=1.5pt, shorten <=1.5pt},
    xlink/.style={->, dashed, line width=0.9pt, redDark!70,
                  shorten >=1.5pt, shorten <=1.5pt},
    hop/.style={preaction={draw, line width=3.4pt, white}},
    junction/.style={circle, fill=black, inner sep=0pt, minimum size=3.4pt},
    %
    % --- boxes: every one carries generous inner sep so text never
    %     touches its own border ---
    dnode/.style={rounded corners=5pt, align=center, font=\small,
                  inner xsep=5pt, inner ysep=4pt, minimum height=9mm},
    io/.style     ={dnode, draw=grayDark,      fill=neutral},
    proc/.style   ={dnode, draw=blue!60!black, fill=attention},
    latentn/.style={dnode, draw=orange!65!black, fill=latent},
    learn/.style  ={dnode, draw=redDark,       fill=flow},
    attn/.style   ={dnode, draw=blue!60!black, fill=attention},
    head/.style   ={dnode, draw=greenDark,     fill=detail},
    cond/.style   ={dnode, draw=purpleDark,    fill=vq},
    seqn/.style   ={dnode, draw=blue!55!black, fill=decoder},
    frozen/.style ={dnode, draw=blue!55!black, fill=decoder},
    %
    % --- tree nodes ---
    troot/.style={dnode, draw=grayDark, fill=neutral, font=\small\bfseries,
                  text width=34mm},
    tcat/.style ={dnode, draw=blue!60!black, fill=attention,
                  font=\small\bfseries, text width=30mm},
    tleaf/.style={dnode, draw=grayDark!70, fill=white, font=\scriptsize,
                  text width=34mm, minimum height=8mm},
    tmine/.style={dnode, draw=redDark, line width=1pt, fill=flow,
                  font=\scriptsize, text width=34mm, minimum height=8mm},
    %
    % --- grouping panels ---
    panelA/.style={rounded corners=9pt, draw=blue!45,       dashed,
                   line width=1pt, fill=panelBlue},
    panelB/.style={rounded corners=9pt, draw=green!45!black, dashed,
                   line width=1pt, fill=panelGreen},
    panelC/.style={rounded corners=9pt, draw=redDark!55,     dashed,
                   line width=1pt, fill=panelWarm},
    %
    % --- annotation ---
    note/.style={font=\scriptsize, align=center, text=grayDark},
    tag/.style={font=\scriptsize, align=center, fill=white, inner sep=1.5pt},
    %
    % --- legend ---
    chip/.style={rounded corners=3pt, minimum width=.65cm,
                 minimum height=.3cm, inner sep=0pt},
    chiptext/.style={anchor=west, font=\scriptsize},
    legendbox/.style={rounded corners=7pt, draw=black!12, fill=white},
}

% Numbered section badge, used to title a panel.
% #1 stroke/fill colour, #2 node name, #3 the number, #4 anchor node,
% #5 the title text.
\newcommand{\panelbadge}[5]{%
  \node[circle, draw=#1, fill=#1, text=white, minimum size=.62cm,
        font=\small\bfseries] (#2) at ($(#4.north west)+(0.72,-0.58)$) {#3};
  \node[anchor=west, font=\bfseries, text=#1, right=2mm of #2] {#5};
}

% A legend swatch chained to the right of the previous one.
% #1 draw colour, #2 fill colour, #3 node name, #4 previous text node,
% #5 label
\newcommand{\legendchip}[5]{%
  \node[chip, draw=#1, fill=#2, right=8mm of #4] (#3) {};
  \node[chiptext, right=2mm of #3] (#3t) {#5};
}

%   \begin{document} \begin{tikzpicture} ... \end{tikzpicture} \end{document}
%
% The palette is the one introduced by diagram-orbis.tex. That file
% still carries its own inline copy of these definitions; if you ever
% touch the palette, change it in both places or switch orbis over to
% % ============================================================
% Shared preamble for every standalone diagram in this folder.
%
% Usage, from a file in this same directory:
%
%   \documentclass[tikz,border=8pt]{standalone}
%   \input{diagramstyle}
%   \begin{document} \begin{tikzpicture} ... \end{tikzpicture} \end{document}
%
% The palette is the one introduced by diagram-orbis.tex. That file
% still carries its own inline copy of these definitions; if you ever
% touch the palette, change it in both places or switch orbis over to
% \input{diagramstyle} as well.
% ============================================================

\usepackage{amsmath,amssymb}
\usepackage{tikz}
\usetikzlibrary{arrows.meta, positioning, calc, fit, backgrounds,
                shapes.geometric, matrix, decorations.pathreplacing}

% ---------- fills -------------------------------------------
% Panel fills are OPAQUE pre-blended colours rather than transparent
% ones, so a line can be given a halo in exactly the panel colour.
% That is what turns an unavoidable crossing into a clean "hop".
\definecolor{panelBlue}{RGB}{248,251,255}
\definecolor{panelGreen}{RGB}{249,253,249}
\definecolor{panelWarm}{RGB}{255,251,247}

\definecolor{semantic}{RGB}{229,224,250}
\definecolor{detail}{RGB}{218,246,228}
\definecolor{latent}{RGB}{255,244,210}
\definecolor{vq}{RGB}{244,232,255}
\definecolor{flow}{RGB}{255,231,222}
\definecolor{attention}{RGB}{224,239,255}
\definecolor{decoder}{RGB}{231,236,255}
\definecolor{neutral}{RGB}{243,244,246}

% ---------- strokes -----------------------------------------
\definecolor{blueDark}{RGB}{25,73,115}
\definecolor{greenDark}{RGB}{35,120,78}
\definecolor{purpleDark}{RGB}{105,70,145}
\definecolor{redDark}{RGB}{145,65,55}
\definecolor{grayDark}{RGB}{78,82,90}
\definecolor{orangeDark}{RGB}{165,95,40}

% ---------- styles ------------------------------------------
\tikzset{
    >=Latex,
    % --- connectors ---
    arrow/.style={->, line width=0.8pt, shorten >=1.5pt, shorten <=1.5pt},
    thinarrow/.style={->, line width=0.6pt, shorten >=1.2pt, shorten <=1.2pt},
    plainline/.style={line width=0.8pt},
    dashedarrow/.style={->, dashed, line width=0.8pt,
                        shorten >=1.5pt, shorten <=1.5pt},
    xlink/.style={->, dashed, line width=0.9pt, redDark!70,
                  shorten >=1.5pt, shorten <=1.5pt},
    hop/.style={preaction={draw, line width=3.4pt, white}},
    junction/.style={circle, fill=black, inner sep=0pt, minimum size=3.4pt},
    %
    % --- boxes: every one carries generous inner sep so text never
    %     touches its own border ---
    dnode/.style={rounded corners=5pt, align=center, font=\small,
                  inner xsep=5pt, inner ysep=4pt, minimum height=9mm},
    io/.style     ={dnode, draw=grayDark,      fill=neutral},
    proc/.style   ={dnode, draw=blue!60!black, fill=attention},
    latentn/.style={dnode, draw=orange!65!black, fill=latent},
    learn/.style  ={dnode, draw=redDark,       fill=flow},
    attn/.style   ={dnode, draw=blue!60!black, fill=attention},
    head/.style   ={dnode, draw=greenDark,     fill=detail},
    cond/.style   ={dnode, draw=purpleDark,    fill=vq},
    seqn/.style   ={dnode, draw=blue!55!black, fill=decoder},
    frozen/.style ={dnode, draw=blue!55!black, fill=decoder},
    %
    % --- tree nodes ---
    troot/.style={dnode, draw=grayDark, fill=neutral, font=\small\bfseries,
                  text width=34mm},
    tcat/.style ={dnode, draw=blue!60!black, fill=attention,
                  font=\small\bfseries, text width=30mm},
    tleaf/.style={dnode, draw=grayDark!70, fill=white, font=\scriptsize,
                  text width=34mm, minimum height=8mm},
    tmine/.style={dnode, draw=redDark, line width=1pt, fill=flow,
                  font=\scriptsize, text width=34mm, minimum height=8mm},
    %
    % --- grouping panels ---
    panelA/.style={rounded corners=9pt, draw=blue!45,       dashed,
                   line width=1pt, fill=panelBlue},
    panelB/.style={rounded corners=9pt, draw=green!45!black, dashed,
                   line width=1pt, fill=panelGreen},
    panelC/.style={rounded corners=9pt, draw=redDark!55,     dashed,
                   line width=1pt, fill=panelWarm},
    %
    % --- annotation ---
    note/.style={font=\scriptsize, align=center, text=grayDark},
    tag/.style={font=\scriptsize, align=center, fill=white, inner sep=1.5pt},
    %
    % --- legend ---
    chip/.style={rounded corners=3pt, minimum width=.65cm,
                 minimum height=.3cm, inner sep=0pt},
    chiptext/.style={anchor=west, font=\scriptsize},
    legendbox/.style={rounded corners=7pt, draw=black!12, fill=white},
}

% Numbered section badge, used to title a panel.
% #1 stroke/fill colour, #2 node name, #3 the number, #4 anchor node,
% #5 the title text.
\newcommand{\panelbadge}[5]{%
  \node[circle, draw=#1, fill=#1, text=white, minimum size=.62cm,
        font=\small\bfseries] (#2) at ($(#4.north west)+(0.72,-0.58)$) {#3};
  \node[anchor=west, font=\bfseries, text=#1, right=2mm of #2] {#5};
}

% A legend swatch chained to the right of the previous one.
% #1 draw colour, #2 fill colour, #3 node name, #4 previous text node,
% #5 label
\newcommand{\legendchip}[5]{%
  \node[chip, draw=#1, fill=#2, right=8mm of #4] (#3) {};
  \node[chiptext, right=2mm of #3] (#3t) {#5};
}
 as well.
% ============================================================

\usepackage{amsmath,amssymb}
\usepackage{tikz}
\usetikzlibrary{arrows.meta, positioning, calc, fit, backgrounds,
                shapes.geometric, matrix, decorations.pathreplacing}

% ---------- fills -------------------------------------------
% Panel fills are OPAQUE pre-blended colours rather than transparent
% ones, so a line can be given a halo in exactly the panel colour.
% That is what turns an unavoidable crossing into a clean "hop".
\definecolor{panelBlue}{RGB}{248,251,255}
\definecolor{panelGreen}{RGB}{249,253,249}
\definecolor{panelWarm}{RGB}{255,251,247}

\definecolor{semantic}{RGB}{229,224,250}
\definecolor{detail}{RGB}{218,246,228}
\definecolor{latent}{RGB}{255,244,210}
\definecolor{vq}{RGB}{244,232,255}
\definecolor{flow}{RGB}{255,231,222}
\definecolor{attention}{RGB}{224,239,255}
\definecolor{decoder}{RGB}{231,236,255}
\definecolor{neutral}{RGB}{243,244,246}

% ---------- strokes -----------------------------------------
\definecolor{blueDark}{RGB}{25,73,115}
\definecolor{greenDark}{RGB}{35,120,78}
\definecolor{purpleDark}{RGB}{105,70,145}
\definecolor{redDark}{RGB}{145,65,55}
\definecolor{grayDark}{RGB}{78,82,90}
\definecolor{orangeDark}{RGB}{165,95,40}

% ---------- styles ------------------------------------------
\tikzset{
    >=Latex,
    % --- connectors ---
    arrow/.style={->, line width=0.8pt, shorten >=1.5pt, shorten <=1.5pt},
    thinarrow/.style={->, line width=0.6pt, shorten >=1.2pt, shorten <=1.2pt},
    plainline/.style={line width=0.8pt},
    dashedarrow/.style={->, dashed, line width=0.8pt,
                        shorten >=1.5pt, shorten <=1.5pt},
    xlink/.style={->, dashed, line width=0.9pt, redDark!70,
                  shorten >=1.5pt, shorten <=1.5pt},
    hop/.style={preaction={draw, line width=3.4pt, white}},
    junction/.style={circle, fill=black, inner sep=0pt, minimum size=3.4pt},
    %
    % --- boxes: every one carries generous inner sep so text never
    %     touches its own border ---
    dnode/.style={rounded corners=5pt, align=center, font=\small,
                  inner xsep=5pt, inner ysep=4pt, minimum height=9mm},
    io/.style     ={dnode, draw=grayDark,      fill=neutral},
    proc/.style   ={dnode, draw=blue!60!black, fill=attention},
    latentn/.style={dnode, draw=orange!65!black, fill=latent},
    learn/.style  ={dnode, draw=redDark,       fill=flow},
    attn/.style   ={dnode, draw=blue!60!black, fill=attention},
    head/.style   ={dnode, draw=greenDark,     fill=detail},
    cond/.style   ={dnode, draw=purpleDark,    fill=vq},
    seqn/.style   ={dnode, draw=blue!55!black, fill=decoder},
    frozen/.style ={dnode, draw=blue!55!black, fill=decoder},
    %
    % --- tree nodes ---
    troot/.style={dnode, draw=grayDark, fill=neutral, font=\small\bfseries,
                  text width=34mm},
    tcat/.style ={dnode, draw=blue!60!black, fill=attention,
                  font=\small\bfseries, text width=30mm},
    tleaf/.style={dnode, draw=grayDark!70, fill=white, font=\scriptsize,
                  text width=34mm, minimum height=8mm},
    tmine/.style={dnode, draw=redDark, line width=1pt, fill=flow,
                  font=\scriptsize, text width=34mm, minimum height=8mm},
    %
    % --- grouping panels ---
    panelA/.style={rounded corners=9pt, draw=blue!45,       dashed,
                   line width=1pt, fill=panelBlue},
    panelB/.style={rounded corners=9pt, draw=green!45!black, dashed,
                   line width=1pt, fill=panelGreen},
    panelC/.style={rounded corners=9pt, draw=redDark!55,     dashed,
                   line width=1pt, fill=panelWarm},
    %
    % --- annotation ---
    note/.style={font=\scriptsize, align=center, text=grayDark},
    tag/.style={font=\scriptsize, align=center, fill=white, inner sep=1.5pt},
    %
    % --- legend ---
    chip/.style={rounded corners=3pt, minimum width=.65cm,
                 minimum height=.3cm, inner sep=0pt},
    chiptext/.style={anchor=west, font=\scriptsize},
    legendbox/.style={rounded corners=7pt, draw=black!12, fill=white},
}

% Numbered section badge, used to title a panel.
% #1 stroke/fill colour, #2 node name, #3 the number, #4 anchor node,
% #5 the title text.
\newcommand{\panelbadge}[5]{%
  \node[circle, draw=#1, fill=#1, text=white, minimum size=.62cm,
        font=\small\bfseries] (#2) at ($(#4.north west)+(0.72,-0.58)$) {#3};
  \node[anchor=west, font=\bfseries, text=#1, right=2mm of #2] {#5};
}

% A legend swatch chained to the right of the previous one.
% #1 draw colour, #2 fill colour, #3 node name, #4 previous text node,
% #5 label
\newcommand{\legendchip}[5]{%
  \node[chip, draw=#1, fill=#2, right=8mm of #4] (#3) {};
  \node[chiptext, right=2mm of #3] (#3t) {#5};
}

%   \begin{document} \begin{tikzpicture} ... \end{tikzpicture} \end{document}
%
% The palette is the one introduced by diagram-orbis.tex. That file
% still carries its own inline copy of these definitions; if you ever
% touch the palette, change it in both places or switch orbis over to
% % ============================================================
% Shared preamble for every standalone diagram in this folder.
%
% Usage, from a file in this same directory:
%
%   \documentclass[tikz,border=8pt]{standalone}
%   % ============================================================
% Shared preamble for every standalone diagram in this folder.
%
% Usage, from a file in this same directory:
%
%   \documentclass[tikz,border=8pt]{standalone}
%   \input{diagramstyle}
%   \begin{document} \begin{tikzpicture} ... \end{tikzpicture} \end{document}
%
% The palette is the one introduced by diagram-orbis.tex. That file
% still carries its own inline copy of these definitions; if you ever
% touch the palette, change it in both places or switch orbis over to
% \input{diagramstyle} as well.
% ============================================================

\usepackage{amsmath,amssymb}
\usepackage{tikz}
\usetikzlibrary{arrows.meta, positioning, calc, fit, backgrounds,
                shapes.geometric, matrix, decorations.pathreplacing}

% ---------- fills -------------------------------------------
% Panel fills are OPAQUE pre-blended colours rather than transparent
% ones, so a line can be given a halo in exactly the panel colour.
% That is what turns an unavoidable crossing into a clean "hop".
\definecolor{panelBlue}{RGB}{248,251,255}
\definecolor{panelGreen}{RGB}{249,253,249}
\definecolor{panelWarm}{RGB}{255,251,247}

\definecolor{semantic}{RGB}{229,224,250}
\definecolor{detail}{RGB}{218,246,228}
\definecolor{latent}{RGB}{255,244,210}
\definecolor{vq}{RGB}{244,232,255}
\definecolor{flow}{RGB}{255,231,222}
\definecolor{attention}{RGB}{224,239,255}
\definecolor{decoder}{RGB}{231,236,255}
\definecolor{neutral}{RGB}{243,244,246}

% ---------- strokes -----------------------------------------
\definecolor{blueDark}{RGB}{25,73,115}
\definecolor{greenDark}{RGB}{35,120,78}
\definecolor{purpleDark}{RGB}{105,70,145}
\definecolor{redDark}{RGB}{145,65,55}
\definecolor{grayDark}{RGB}{78,82,90}
\definecolor{orangeDark}{RGB}{165,95,40}

% ---------- styles ------------------------------------------
\tikzset{
    >=Latex,
    % --- connectors ---
    arrow/.style={->, line width=0.8pt, shorten >=1.5pt, shorten <=1.5pt},
    thinarrow/.style={->, line width=0.6pt, shorten >=1.2pt, shorten <=1.2pt},
    plainline/.style={line width=0.8pt},
    dashedarrow/.style={->, dashed, line width=0.8pt,
                        shorten >=1.5pt, shorten <=1.5pt},
    xlink/.style={->, dashed, line width=0.9pt, redDark!70,
                  shorten >=1.5pt, shorten <=1.5pt},
    hop/.style={preaction={draw, line width=3.4pt, white}},
    junction/.style={circle, fill=black, inner sep=0pt, minimum size=3.4pt},
    %
    % --- boxes: every one carries generous inner sep so text never
    %     touches its own border ---
    dnode/.style={rounded corners=5pt, align=center, font=\small,
                  inner xsep=5pt, inner ysep=4pt, minimum height=9mm},
    io/.style     ={dnode, draw=grayDark,      fill=neutral},
    proc/.style   ={dnode, draw=blue!60!black, fill=attention},
    latentn/.style={dnode, draw=orange!65!black, fill=latent},
    learn/.style  ={dnode, draw=redDark,       fill=flow},
    attn/.style   ={dnode, draw=blue!60!black, fill=attention},
    head/.style   ={dnode, draw=greenDark,     fill=detail},
    cond/.style   ={dnode, draw=purpleDark,    fill=vq},
    seqn/.style   ={dnode, draw=blue!55!black, fill=decoder},
    frozen/.style ={dnode, draw=blue!55!black, fill=decoder},
    %
    % --- tree nodes ---
    troot/.style={dnode, draw=grayDark, fill=neutral, font=\small\bfseries,
                  text width=34mm},
    tcat/.style ={dnode, draw=blue!60!black, fill=attention,
                  font=\small\bfseries, text width=30mm},
    tleaf/.style={dnode, draw=grayDark!70, fill=white, font=\scriptsize,
                  text width=34mm, minimum height=8mm},
    tmine/.style={dnode, draw=redDark, line width=1pt, fill=flow,
                  font=\scriptsize, text width=34mm, minimum height=8mm},
    %
    % --- grouping panels ---
    panelA/.style={rounded corners=9pt, draw=blue!45,       dashed,
                   line width=1pt, fill=panelBlue},
    panelB/.style={rounded corners=9pt, draw=green!45!black, dashed,
                   line width=1pt, fill=panelGreen},
    panelC/.style={rounded corners=9pt, draw=redDark!55,     dashed,
                   line width=1pt, fill=panelWarm},
    %
    % --- annotation ---
    note/.style={font=\scriptsize, align=center, text=grayDark},
    tag/.style={font=\scriptsize, align=center, fill=white, inner sep=1.5pt},
    %
    % --- legend ---
    chip/.style={rounded corners=3pt, minimum width=.65cm,
                 minimum height=.3cm, inner sep=0pt},
    chiptext/.style={anchor=west, font=\scriptsize},
    legendbox/.style={rounded corners=7pt, draw=black!12, fill=white},
}

% Numbered section badge, used to title a panel.
% #1 stroke/fill colour, #2 node name, #3 the number, #4 anchor node,
% #5 the title text.
\newcommand{\panelbadge}[5]{%
  \node[circle, draw=#1, fill=#1, text=white, minimum size=.62cm,
        font=\small\bfseries] (#2) at ($(#4.north west)+(0.72,-0.58)$) {#3};
  \node[anchor=west, font=\bfseries, text=#1, right=2mm of #2] {#5};
}

% A legend swatch chained to the right of the previous one.
% #1 draw colour, #2 fill colour, #3 node name, #4 previous text node,
% #5 label
\newcommand{\legendchip}[5]{%
  \node[chip, draw=#1, fill=#2, right=8mm of #4] (#3) {};
  \node[chiptext, right=2mm of #3] (#3t) {#5};
}

%   \begin{document} \begin{tikzpicture} ... \end{tikzpicture} \end{document}
%
% The palette is the one introduced by diagram-orbis.tex. That file
% still carries its own inline copy of these definitions; if you ever
% touch the palette, change it in both places or switch orbis over to
% % ============================================================
% Shared preamble for every standalone diagram in this folder.
%
% Usage, from a file in this same directory:
%
%   \documentclass[tikz,border=8pt]{standalone}
%   \input{diagramstyle}
%   \begin{document} \begin{tikzpicture} ... \end{tikzpicture} \end{document}
%
% The palette is the one introduced by diagram-orbis.tex. That file
% still carries its own inline copy of these definitions; if you ever
% touch the palette, change it in both places or switch orbis over to
% \input{diagramstyle} as well.
% ============================================================

\usepackage{amsmath,amssymb}
\usepackage{tikz}
\usetikzlibrary{arrows.meta, positioning, calc, fit, backgrounds,
                shapes.geometric, matrix, decorations.pathreplacing}

% ---------- fills -------------------------------------------
% Panel fills are OPAQUE pre-blended colours rather than transparent
% ones, so a line can be given a halo in exactly the panel colour.
% That is what turns an unavoidable crossing into a clean "hop".
\definecolor{panelBlue}{RGB}{248,251,255}
\definecolor{panelGreen}{RGB}{249,253,249}
\definecolor{panelWarm}{RGB}{255,251,247}

\definecolor{semantic}{RGB}{229,224,250}
\definecolor{detail}{RGB}{218,246,228}
\definecolor{latent}{RGB}{255,244,210}
\definecolor{vq}{RGB}{244,232,255}
\definecolor{flow}{RGB}{255,231,222}
\definecolor{attention}{RGB}{224,239,255}
\definecolor{decoder}{RGB}{231,236,255}
\definecolor{neutral}{RGB}{243,244,246}

% ---------- strokes -----------------------------------------
\definecolor{blueDark}{RGB}{25,73,115}
\definecolor{greenDark}{RGB}{35,120,78}
\definecolor{purpleDark}{RGB}{105,70,145}
\definecolor{redDark}{RGB}{145,65,55}
\definecolor{grayDark}{RGB}{78,82,90}
\definecolor{orangeDark}{RGB}{165,95,40}

% ---------- styles ------------------------------------------
\tikzset{
    >=Latex,
    % --- connectors ---
    arrow/.style={->, line width=0.8pt, shorten >=1.5pt, shorten <=1.5pt},
    thinarrow/.style={->, line width=0.6pt, shorten >=1.2pt, shorten <=1.2pt},
    plainline/.style={line width=0.8pt},
    dashedarrow/.style={->, dashed, line width=0.8pt,
                        shorten >=1.5pt, shorten <=1.5pt},
    xlink/.style={->, dashed, line width=0.9pt, redDark!70,
                  shorten >=1.5pt, shorten <=1.5pt},
    hop/.style={preaction={draw, line width=3.4pt, white}},
    junction/.style={circle, fill=black, inner sep=0pt, minimum size=3.4pt},
    %
    % --- boxes: every one carries generous inner sep so text never
    %     touches its own border ---
    dnode/.style={rounded corners=5pt, align=center, font=\small,
                  inner xsep=5pt, inner ysep=4pt, minimum height=9mm},
    io/.style     ={dnode, draw=grayDark,      fill=neutral},
    proc/.style   ={dnode, draw=blue!60!black, fill=attention},
    latentn/.style={dnode, draw=orange!65!black, fill=latent},
    learn/.style  ={dnode, draw=redDark,       fill=flow},
    attn/.style   ={dnode, draw=blue!60!black, fill=attention},
    head/.style   ={dnode, draw=greenDark,     fill=detail},
    cond/.style   ={dnode, draw=purpleDark,    fill=vq},
    seqn/.style   ={dnode, draw=blue!55!black, fill=decoder},
    frozen/.style ={dnode, draw=blue!55!black, fill=decoder},
    %
    % --- tree nodes ---
    troot/.style={dnode, draw=grayDark, fill=neutral, font=\small\bfseries,
                  text width=34mm},
    tcat/.style ={dnode, draw=blue!60!black, fill=attention,
                  font=\small\bfseries, text width=30mm},
    tleaf/.style={dnode, draw=grayDark!70, fill=white, font=\scriptsize,
                  text width=34mm, minimum height=8mm},
    tmine/.style={dnode, draw=redDark, line width=1pt, fill=flow,
                  font=\scriptsize, text width=34mm, minimum height=8mm},
    %
    % --- grouping panels ---
    panelA/.style={rounded corners=9pt, draw=blue!45,       dashed,
                   line width=1pt, fill=panelBlue},
    panelB/.style={rounded corners=9pt, draw=green!45!black, dashed,
                   line width=1pt, fill=panelGreen},
    panelC/.style={rounded corners=9pt, draw=redDark!55,     dashed,
                   line width=1pt, fill=panelWarm},
    %
    % --- annotation ---
    note/.style={font=\scriptsize, align=center, text=grayDark},
    tag/.style={font=\scriptsize, align=center, fill=white, inner sep=1.5pt},
    %
    % --- legend ---
    chip/.style={rounded corners=3pt, minimum width=.65cm,
                 minimum height=.3cm, inner sep=0pt},
    chiptext/.style={anchor=west, font=\scriptsize},
    legendbox/.style={rounded corners=7pt, draw=black!12, fill=white},
}

% Numbered section badge, used to title a panel.
% #1 stroke/fill colour, #2 node name, #3 the number, #4 anchor node,
% #5 the title text.
\newcommand{\panelbadge}[5]{%
  \node[circle, draw=#1, fill=#1, text=white, minimum size=.62cm,
        font=\small\bfseries] (#2) at ($(#4.north west)+(0.72,-0.58)$) {#3};
  \node[anchor=west, font=\bfseries, text=#1, right=2mm of #2] {#5};
}

% A legend swatch chained to the right of the previous one.
% #1 draw colour, #2 fill colour, #3 node name, #4 previous text node,
% #5 label
\newcommand{\legendchip}[5]{%
  \node[chip, draw=#1, fill=#2, right=8mm of #4] (#3) {};
  \node[chiptext, right=2mm of #3] (#3t) {#5};
}
 as well.
% ============================================================

\usepackage{amsmath,amssymb}
\usepackage{tikz}
\usetikzlibrary{arrows.meta, positioning, calc, fit, backgrounds,
                shapes.geometric, matrix, decorations.pathreplacing}

% ---------- fills -------------------------------------------
% Panel fills are OPAQUE pre-blended colours rather than transparent
% ones, so a line can be given a halo in exactly the panel colour.
% That is what turns an unavoidable crossing into a clean "hop".
\definecolor{panelBlue}{RGB}{248,251,255}
\definecolor{panelGreen}{RGB}{249,253,249}
\definecolor{panelWarm}{RGB}{255,251,247}

\definecolor{semantic}{RGB}{229,224,250}
\definecolor{detail}{RGB}{218,246,228}
\definecolor{latent}{RGB}{255,244,210}
\definecolor{vq}{RGB}{244,232,255}
\definecolor{flow}{RGB}{255,231,222}
\definecolor{attention}{RGB}{224,239,255}
\definecolor{decoder}{RGB}{231,236,255}
\definecolor{neutral}{RGB}{243,244,246}

% ---------- strokes -----------------------------------------
\definecolor{blueDark}{RGB}{25,73,115}
\definecolor{greenDark}{RGB}{35,120,78}
\definecolor{purpleDark}{RGB}{105,70,145}
\definecolor{redDark}{RGB}{145,65,55}
\definecolor{grayDark}{RGB}{78,82,90}
\definecolor{orangeDark}{RGB}{165,95,40}

% ---------- styles ------------------------------------------
\tikzset{
    >=Latex,
    % --- connectors ---
    arrow/.style={->, line width=0.8pt, shorten >=1.5pt, shorten <=1.5pt},
    thinarrow/.style={->, line width=0.6pt, shorten >=1.2pt, shorten <=1.2pt},
    plainline/.style={line width=0.8pt},
    dashedarrow/.style={->, dashed, line width=0.8pt,
                        shorten >=1.5pt, shorten <=1.5pt},
    xlink/.style={->, dashed, line width=0.9pt, redDark!70,
                  shorten >=1.5pt, shorten <=1.5pt},
    hop/.style={preaction={draw, line width=3.4pt, white}},
    junction/.style={circle, fill=black, inner sep=0pt, minimum size=3.4pt},
    %
    % --- boxes: every one carries generous inner sep so text never
    %     touches its own border ---
    dnode/.style={rounded corners=5pt, align=center, font=\small,
                  inner xsep=5pt, inner ysep=4pt, minimum height=9mm},
    io/.style     ={dnode, draw=grayDark,      fill=neutral},
    proc/.style   ={dnode, draw=blue!60!black, fill=attention},
    latentn/.style={dnode, draw=orange!65!black, fill=latent},
    learn/.style  ={dnode, draw=redDark,       fill=flow},
    attn/.style   ={dnode, draw=blue!60!black, fill=attention},
    head/.style   ={dnode, draw=greenDark,     fill=detail},
    cond/.style   ={dnode, draw=purpleDark,    fill=vq},
    seqn/.style   ={dnode, draw=blue!55!black, fill=decoder},
    frozen/.style ={dnode, draw=blue!55!black, fill=decoder},
    %
    % --- tree nodes ---
    troot/.style={dnode, draw=grayDark, fill=neutral, font=\small\bfseries,
                  text width=34mm},
    tcat/.style ={dnode, draw=blue!60!black, fill=attention,
                  font=\small\bfseries, text width=30mm},
    tleaf/.style={dnode, draw=grayDark!70, fill=white, font=\scriptsize,
                  text width=34mm, minimum height=8mm},
    tmine/.style={dnode, draw=redDark, line width=1pt, fill=flow,
                  font=\scriptsize, text width=34mm, minimum height=8mm},
    %
    % --- grouping panels ---
    panelA/.style={rounded corners=9pt, draw=blue!45,       dashed,
                   line width=1pt, fill=panelBlue},
    panelB/.style={rounded corners=9pt, draw=green!45!black, dashed,
                   line width=1pt, fill=panelGreen},
    panelC/.style={rounded corners=9pt, draw=redDark!55,     dashed,
                   line width=1pt, fill=panelWarm},
    %
    % --- annotation ---
    note/.style={font=\scriptsize, align=center, text=grayDark},
    tag/.style={font=\scriptsize, align=center, fill=white, inner sep=1.5pt},
    %
    % --- legend ---
    chip/.style={rounded corners=3pt, minimum width=.65cm,
                 minimum height=.3cm, inner sep=0pt},
    chiptext/.style={anchor=west, font=\scriptsize},
    legendbox/.style={rounded corners=7pt, draw=black!12, fill=white},
}

% Numbered section badge, used to title a panel.
% #1 stroke/fill colour, #2 node name, #3 the number, #4 anchor node,
% #5 the title text.
\newcommand{\panelbadge}[5]{%
  \node[circle, draw=#1, fill=#1, text=white, minimum size=.62cm,
        font=\small\bfseries] (#2) at ($(#4.north west)+(0.72,-0.58)$) {#3};
  \node[anchor=west, font=\bfseries, text=#1, right=2mm of #2] {#5};
}

% A legend swatch chained to the right of the previous one.
% #1 draw colour, #2 fill colour, #3 node name, #4 previous text node,
% #5 label
\newcommand{\legendchip}[5]{%
  \node[chip, draw=#1, fill=#2, right=8mm of #4] (#3) {};
  \node[chiptext, right=2mm of #3] (#3t) {#5};
}
 as well.
% ============================================================

\usepackage{amsmath,amssymb}
\usepackage{tikz}
\usetikzlibrary{arrows.meta, positioning, calc, fit, backgrounds,
                shapes.geometric, matrix, decorations.pathreplacing}

% ---------- fills -------------------------------------------
% Panel fills are OPAQUE pre-blended colours rather than transparent
% ones, so a line can be given a halo in exactly the panel colour.
% That is what turns an unavoidable crossing into a clean "hop".
\definecolor{panelBlue}{RGB}{248,251,255}
\definecolor{panelGreen}{RGB}{249,253,249}
\definecolor{panelWarm}{RGB}{255,251,247}

\definecolor{semantic}{RGB}{229,224,250}
\definecolor{detail}{RGB}{218,246,228}
\definecolor{latent}{RGB}{255,244,210}
\definecolor{vq}{RGB}{244,232,255}
\definecolor{flow}{RGB}{255,231,222}
\definecolor{attention}{RGB}{224,239,255}
\definecolor{decoder}{RGB}{231,236,255}
\definecolor{neutral}{RGB}{243,244,246}

% ---------- strokes -----------------------------------------
\definecolor{blueDark}{RGB}{25,73,115}
\definecolor{greenDark}{RGB}{35,120,78}
\definecolor{purpleDark}{RGB}{105,70,145}
\definecolor{redDark}{RGB}{145,65,55}
\definecolor{grayDark}{RGB}{78,82,90}
\definecolor{orangeDark}{RGB}{165,95,40}

% ---------- styles ------------------------------------------
\tikzset{
    >=Latex,
    % --- connectors ---
    arrow/.style={->, line width=0.8pt, shorten >=1.5pt, shorten <=1.5pt},
    thinarrow/.style={->, line width=0.6pt, shorten >=1.2pt, shorten <=1.2pt},
    plainline/.style={line width=0.8pt},
    dashedarrow/.style={->, dashed, line width=0.8pt,
                        shorten >=1.5pt, shorten <=1.5pt},
    xlink/.style={->, dashed, line width=0.9pt, redDark!70,
                  shorten >=1.5pt, shorten <=1.5pt},
    hop/.style={preaction={draw, line width=3.4pt, white}},
    junction/.style={circle, fill=black, inner sep=0pt, minimum size=3.4pt},
    %
    % --- boxes: every one carries generous inner sep so text never
    %     touches its own border ---
    dnode/.style={rounded corners=5pt, align=center, font=\small,
                  inner xsep=5pt, inner ysep=4pt, minimum height=9mm},
    io/.style     ={dnode, draw=grayDark,      fill=neutral},
    proc/.style   ={dnode, draw=blue!60!black, fill=attention},
    latentn/.style={dnode, draw=orange!65!black, fill=latent},
    learn/.style  ={dnode, draw=redDark,       fill=flow},
    attn/.style   ={dnode, draw=blue!60!black, fill=attention},
    head/.style   ={dnode, draw=greenDark,     fill=detail},
    cond/.style   ={dnode, draw=purpleDark,    fill=vq},
    seqn/.style   ={dnode, draw=blue!55!black, fill=decoder},
    frozen/.style ={dnode, draw=blue!55!black, fill=decoder},
    %
    % --- tree nodes ---
    troot/.style={dnode, draw=grayDark, fill=neutral, font=\small\bfseries,
                  text width=34mm},
    tcat/.style ={dnode, draw=blue!60!black, fill=attention,
                  font=\small\bfseries, text width=30mm},
    tleaf/.style={dnode, draw=grayDark!70, fill=white, font=\scriptsize,
                  text width=34mm, minimum height=8mm},
    tmine/.style={dnode, draw=redDark, line width=1pt, fill=flow,
                  font=\scriptsize, text width=34mm, minimum height=8mm},
    %
    % --- grouping panels ---
    panelA/.style={rounded corners=9pt, draw=blue!45,       dashed,
                   line width=1pt, fill=panelBlue},
    panelB/.style={rounded corners=9pt, draw=green!45!black, dashed,
                   line width=1pt, fill=panelGreen},
    panelC/.style={rounded corners=9pt, draw=redDark!55,     dashed,
                   line width=1pt, fill=panelWarm},
    %
    % --- annotation ---
    note/.style={font=\scriptsize, align=center, text=grayDark},
    tag/.style={font=\scriptsize, align=center, fill=white, inner sep=1.5pt},
    %
    % --- legend ---
    chip/.style={rounded corners=3pt, minimum width=.65cm,
                 minimum height=.3cm, inner sep=0pt},
    chiptext/.style={anchor=west, font=\scriptsize},
    legendbox/.style={rounded corners=7pt, draw=black!12, fill=white},
}

% Numbered section badge, used to title a panel.
% #1 stroke/fill colour, #2 node name, #3 the number, #4 anchor node,
% #5 the title text.
\newcommand{\panelbadge}[5]{%
  \node[circle, draw=#1, fill=#1, text=white, minimum size=.62cm,
        font=\small\bfseries] (#2) at ($(#4.north west)+(0.72,-0.58)$) {#3};
  \node[anchor=west, font=\bfseries, text=#1, right=2mm of #2] {#5};
}

% A legend swatch chained to the right of the previous one.
% #1 draw colour, #2 fill colour, #3 node name, #4 previous text node,
% #5 label
\newcommand{\legendchip}[5]{%
  \node[chip, draw=#1, fill=#2, right=8mm of #4] (#3) {};
  \node[chiptext, right=2mm of #3] (#3t) {#5};
}

%   \begin{document} \begin{tikzpicture} ... \end{tikzpicture} \end{document}
%
% The palette is the one introduced by diagram-orbis.tex. That file
% still carries its own inline copy of these definitions; if you ever
% touch the palette, change it in both places or switch orbis over to
% % ============================================================
% Shared preamble for every standalone diagram in this folder.
%
% Usage, from a file in this same directory:
%
%   \documentclass[tikz,border=8pt]{standalone}
%   % ============================================================
% Shared preamble for every standalone diagram in this folder.
%
% Usage, from a file in this same directory:
%
%   \documentclass[tikz,border=8pt]{standalone}
%   % ============================================================
% Shared preamble for every standalone diagram in this folder.
%
% Usage, from a file in this same directory:
%
%   \documentclass[tikz,border=8pt]{standalone}
%   \input{diagramstyle}
%   \begin{document} \begin{tikzpicture} ... \end{tikzpicture} \end{document}
%
% The palette is the one introduced by diagram-orbis.tex. That file
% still carries its own inline copy of these definitions; if you ever
% touch the palette, change it in both places or switch orbis over to
% \input{diagramstyle} as well.
% ============================================================

\usepackage{amsmath,amssymb}
\usepackage{tikz}
\usetikzlibrary{arrows.meta, positioning, calc, fit, backgrounds,
                shapes.geometric, matrix, decorations.pathreplacing}

% ---------- fills -------------------------------------------
% Panel fills are OPAQUE pre-blended colours rather than transparent
% ones, so a line can be given a halo in exactly the panel colour.
% That is what turns an unavoidable crossing into a clean "hop".
\definecolor{panelBlue}{RGB}{248,251,255}
\definecolor{panelGreen}{RGB}{249,253,249}
\definecolor{panelWarm}{RGB}{255,251,247}

\definecolor{semantic}{RGB}{229,224,250}
\definecolor{detail}{RGB}{218,246,228}
\definecolor{latent}{RGB}{255,244,210}
\definecolor{vq}{RGB}{244,232,255}
\definecolor{flow}{RGB}{255,231,222}
\definecolor{attention}{RGB}{224,239,255}
\definecolor{decoder}{RGB}{231,236,255}
\definecolor{neutral}{RGB}{243,244,246}

% ---------- strokes -----------------------------------------
\definecolor{blueDark}{RGB}{25,73,115}
\definecolor{greenDark}{RGB}{35,120,78}
\definecolor{purpleDark}{RGB}{105,70,145}
\definecolor{redDark}{RGB}{145,65,55}
\definecolor{grayDark}{RGB}{78,82,90}
\definecolor{orangeDark}{RGB}{165,95,40}

% ---------- styles ------------------------------------------
\tikzset{
    >=Latex,
    % --- connectors ---
    arrow/.style={->, line width=0.8pt, shorten >=1.5pt, shorten <=1.5pt},
    thinarrow/.style={->, line width=0.6pt, shorten >=1.2pt, shorten <=1.2pt},
    plainline/.style={line width=0.8pt},
    dashedarrow/.style={->, dashed, line width=0.8pt,
                        shorten >=1.5pt, shorten <=1.5pt},
    xlink/.style={->, dashed, line width=0.9pt, redDark!70,
                  shorten >=1.5pt, shorten <=1.5pt},
    hop/.style={preaction={draw, line width=3.4pt, white}},
    junction/.style={circle, fill=black, inner sep=0pt, minimum size=3.4pt},
    %
    % --- boxes: every one carries generous inner sep so text never
    %     touches its own border ---
    dnode/.style={rounded corners=5pt, align=center, font=\small,
                  inner xsep=5pt, inner ysep=4pt, minimum height=9mm},
    io/.style     ={dnode, draw=grayDark,      fill=neutral},
    proc/.style   ={dnode, draw=blue!60!black, fill=attention},
    latentn/.style={dnode, draw=orange!65!black, fill=latent},
    learn/.style  ={dnode, draw=redDark,       fill=flow},
    attn/.style   ={dnode, draw=blue!60!black, fill=attention},
    head/.style   ={dnode, draw=greenDark,     fill=detail},
    cond/.style   ={dnode, draw=purpleDark,    fill=vq},
    seqn/.style   ={dnode, draw=blue!55!black, fill=decoder},
    frozen/.style ={dnode, draw=blue!55!black, fill=decoder},
    %
    % --- tree nodes ---
    troot/.style={dnode, draw=grayDark, fill=neutral, font=\small\bfseries,
                  text width=34mm},
    tcat/.style ={dnode, draw=blue!60!black, fill=attention,
                  font=\small\bfseries, text width=30mm},
    tleaf/.style={dnode, draw=grayDark!70, fill=white, font=\scriptsize,
                  text width=34mm, minimum height=8mm},
    tmine/.style={dnode, draw=redDark, line width=1pt, fill=flow,
                  font=\scriptsize, text width=34mm, minimum height=8mm},
    %
    % --- grouping panels ---
    panelA/.style={rounded corners=9pt, draw=blue!45,       dashed,
                   line width=1pt, fill=panelBlue},
    panelB/.style={rounded corners=9pt, draw=green!45!black, dashed,
                   line width=1pt, fill=panelGreen},
    panelC/.style={rounded corners=9pt, draw=redDark!55,     dashed,
                   line width=1pt, fill=panelWarm},
    %
    % --- annotation ---
    note/.style={font=\scriptsize, align=center, text=grayDark},
    tag/.style={font=\scriptsize, align=center, fill=white, inner sep=1.5pt},
    %
    % --- legend ---
    chip/.style={rounded corners=3pt, minimum width=.65cm,
                 minimum height=.3cm, inner sep=0pt},
    chiptext/.style={anchor=west, font=\scriptsize},
    legendbox/.style={rounded corners=7pt, draw=black!12, fill=white},
}

% Numbered section badge, used to title a panel.
% #1 stroke/fill colour, #2 node name, #3 the number, #4 anchor node,
% #5 the title text.
\newcommand{\panelbadge}[5]{%
  \node[circle, draw=#1, fill=#1, text=white, minimum size=.62cm,
        font=\small\bfseries] (#2) at ($(#4.north west)+(0.72,-0.58)$) {#3};
  \node[anchor=west, font=\bfseries, text=#1, right=2mm of #2] {#5};
}

% A legend swatch chained to the right of the previous one.
% #1 draw colour, #2 fill colour, #3 node name, #4 previous text node,
% #5 label
\newcommand{\legendchip}[5]{%
  \node[chip, draw=#1, fill=#2, right=8mm of #4] (#3) {};
  \node[chiptext, right=2mm of #3] (#3t) {#5};
}

%   \begin{document} \begin{tikzpicture} ... \end{tikzpicture} \end{document}
%
% The palette is the one introduced by diagram-orbis.tex. That file
% still carries its own inline copy of these definitions; if you ever
% touch the palette, change it in both places or switch orbis over to
% % ============================================================
% Shared preamble for every standalone diagram in this folder.
%
% Usage, from a file in this same directory:
%
%   \documentclass[tikz,border=8pt]{standalone}
%   \input{diagramstyle}
%   \begin{document} \begin{tikzpicture} ... \end{tikzpicture} \end{document}
%
% The palette is the one introduced by diagram-orbis.tex. That file
% still carries its own inline copy of these definitions; if you ever
% touch the palette, change it in both places or switch orbis over to
% \input{diagramstyle} as well.
% ============================================================

\usepackage{amsmath,amssymb}
\usepackage{tikz}
\usetikzlibrary{arrows.meta, positioning, calc, fit, backgrounds,
                shapes.geometric, matrix, decorations.pathreplacing}

% ---------- fills -------------------------------------------
% Panel fills are OPAQUE pre-blended colours rather than transparent
% ones, so a line can be given a halo in exactly the panel colour.
% That is what turns an unavoidable crossing into a clean "hop".
\definecolor{panelBlue}{RGB}{248,251,255}
\definecolor{panelGreen}{RGB}{249,253,249}
\definecolor{panelWarm}{RGB}{255,251,247}

\definecolor{semantic}{RGB}{229,224,250}
\definecolor{detail}{RGB}{218,246,228}
\definecolor{latent}{RGB}{255,244,210}
\definecolor{vq}{RGB}{244,232,255}
\definecolor{flow}{RGB}{255,231,222}
\definecolor{attention}{RGB}{224,239,255}
\definecolor{decoder}{RGB}{231,236,255}
\definecolor{neutral}{RGB}{243,244,246}

% ---------- strokes -----------------------------------------
\definecolor{blueDark}{RGB}{25,73,115}
\definecolor{greenDark}{RGB}{35,120,78}
\definecolor{purpleDark}{RGB}{105,70,145}
\definecolor{redDark}{RGB}{145,65,55}
\definecolor{grayDark}{RGB}{78,82,90}
\definecolor{orangeDark}{RGB}{165,95,40}

% ---------- styles ------------------------------------------
\tikzset{
    >=Latex,
    % --- connectors ---
    arrow/.style={->, line width=0.8pt, shorten >=1.5pt, shorten <=1.5pt},
    thinarrow/.style={->, line width=0.6pt, shorten >=1.2pt, shorten <=1.2pt},
    plainline/.style={line width=0.8pt},
    dashedarrow/.style={->, dashed, line width=0.8pt,
                        shorten >=1.5pt, shorten <=1.5pt},
    xlink/.style={->, dashed, line width=0.9pt, redDark!70,
                  shorten >=1.5pt, shorten <=1.5pt},
    hop/.style={preaction={draw, line width=3.4pt, white}},
    junction/.style={circle, fill=black, inner sep=0pt, minimum size=3.4pt},
    %
    % --- boxes: every one carries generous inner sep so text never
    %     touches its own border ---
    dnode/.style={rounded corners=5pt, align=center, font=\small,
                  inner xsep=5pt, inner ysep=4pt, minimum height=9mm},
    io/.style     ={dnode, draw=grayDark,      fill=neutral},
    proc/.style   ={dnode, draw=blue!60!black, fill=attention},
    latentn/.style={dnode, draw=orange!65!black, fill=latent},
    learn/.style  ={dnode, draw=redDark,       fill=flow},
    attn/.style   ={dnode, draw=blue!60!black, fill=attention},
    head/.style   ={dnode, draw=greenDark,     fill=detail},
    cond/.style   ={dnode, draw=purpleDark,    fill=vq},
    seqn/.style   ={dnode, draw=blue!55!black, fill=decoder},
    frozen/.style ={dnode, draw=blue!55!black, fill=decoder},
    %
    % --- tree nodes ---
    troot/.style={dnode, draw=grayDark, fill=neutral, font=\small\bfseries,
                  text width=34mm},
    tcat/.style ={dnode, draw=blue!60!black, fill=attention,
                  font=\small\bfseries, text width=30mm},
    tleaf/.style={dnode, draw=grayDark!70, fill=white, font=\scriptsize,
                  text width=34mm, minimum height=8mm},
    tmine/.style={dnode, draw=redDark, line width=1pt, fill=flow,
                  font=\scriptsize, text width=34mm, minimum height=8mm},
    %
    % --- grouping panels ---
    panelA/.style={rounded corners=9pt, draw=blue!45,       dashed,
                   line width=1pt, fill=panelBlue},
    panelB/.style={rounded corners=9pt, draw=green!45!black, dashed,
                   line width=1pt, fill=panelGreen},
    panelC/.style={rounded corners=9pt, draw=redDark!55,     dashed,
                   line width=1pt, fill=panelWarm},
    %
    % --- annotation ---
    note/.style={font=\scriptsize, align=center, text=grayDark},
    tag/.style={font=\scriptsize, align=center, fill=white, inner sep=1.5pt},
    %
    % --- legend ---
    chip/.style={rounded corners=3pt, minimum width=.65cm,
                 minimum height=.3cm, inner sep=0pt},
    chiptext/.style={anchor=west, font=\scriptsize},
    legendbox/.style={rounded corners=7pt, draw=black!12, fill=white},
}

% Numbered section badge, used to title a panel.
% #1 stroke/fill colour, #2 node name, #3 the number, #4 anchor node,
% #5 the title text.
\newcommand{\panelbadge}[5]{%
  \node[circle, draw=#1, fill=#1, text=white, minimum size=.62cm,
        font=\small\bfseries] (#2) at ($(#4.north west)+(0.72,-0.58)$) {#3};
  \node[anchor=west, font=\bfseries, text=#1, right=2mm of #2] {#5};
}

% A legend swatch chained to the right of the previous one.
% #1 draw colour, #2 fill colour, #3 node name, #4 previous text node,
% #5 label
\newcommand{\legendchip}[5]{%
  \node[chip, draw=#1, fill=#2, right=8mm of #4] (#3) {};
  \node[chiptext, right=2mm of #3] (#3t) {#5};
}
 as well.
% ============================================================

\usepackage{amsmath,amssymb}
\usepackage{tikz}
\usetikzlibrary{arrows.meta, positioning, calc, fit, backgrounds,
                shapes.geometric, matrix, decorations.pathreplacing}

% ---------- fills -------------------------------------------
% Panel fills are OPAQUE pre-blended colours rather than transparent
% ones, so a line can be given a halo in exactly the panel colour.
% That is what turns an unavoidable crossing into a clean "hop".
\definecolor{panelBlue}{RGB}{248,251,255}
\definecolor{panelGreen}{RGB}{249,253,249}
\definecolor{panelWarm}{RGB}{255,251,247}

\definecolor{semantic}{RGB}{229,224,250}
\definecolor{detail}{RGB}{218,246,228}
\definecolor{latent}{RGB}{255,244,210}
\definecolor{vq}{RGB}{244,232,255}
\definecolor{flow}{RGB}{255,231,222}
\definecolor{attention}{RGB}{224,239,255}
\definecolor{decoder}{RGB}{231,236,255}
\definecolor{neutral}{RGB}{243,244,246}

% ---------- strokes -----------------------------------------
\definecolor{blueDark}{RGB}{25,73,115}
\definecolor{greenDark}{RGB}{35,120,78}
\definecolor{purpleDark}{RGB}{105,70,145}
\definecolor{redDark}{RGB}{145,65,55}
\definecolor{grayDark}{RGB}{78,82,90}
\definecolor{orangeDark}{RGB}{165,95,40}

% ---------- styles ------------------------------------------
\tikzset{
    >=Latex,
    % --- connectors ---
    arrow/.style={->, line width=0.8pt, shorten >=1.5pt, shorten <=1.5pt},
    thinarrow/.style={->, line width=0.6pt, shorten >=1.2pt, shorten <=1.2pt},
    plainline/.style={line width=0.8pt},
    dashedarrow/.style={->, dashed, line width=0.8pt,
                        shorten >=1.5pt, shorten <=1.5pt},
    xlink/.style={->, dashed, line width=0.9pt, redDark!70,
                  shorten >=1.5pt, shorten <=1.5pt},
    hop/.style={preaction={draw, line width=3.4pt, white}},
    junction/.style={circle, fill=black, inner sep=0pt, minimum size=3.4pt},
    %
    % --- boxes: every one carries generous inner sep so text never
    %     touches its own border ---
    dnode/.style={rounded corners=5pt, align=center, font=\small,
                  inner xsep=5pt, inner ysep=4pt, minimum height=9mm},
    io/.style     ={dnode, draw=grayDark,      fill=neutral},
    proc/.style   ={dnode, draw=blue!60!black, fill=attention},
    latentn/.style={dnode, draw=orange!65!black, fill=latent},
    learn/.style  ={dnode, draw=redDark,       fill=flow},
    attn/.style   ={dnode, draw=blue!60!black, fill=attention},
    head/.style   ={dnode, draw=greenDark,     fill=detail},
    cond/.style   ={dnode, draw=purpleDark,    fill=vq},
    seqn/.style   ={dnode, draw=blue!55!black, fill=decoder},
    frozen/.style ={dnode, draw=blue!55!black, fill=decoder},
    %
    % --- tree nodes ---
    troot/.style={dnode, draw=grayDark, fill=neutral, font=\small\bfseries,
                  text width=34mm},
    tcat/.style ={dnode, draw=blue!60!black, fill=attention,
                  font=\small\bfseries, text width=30mm},
    tleaf/.style={dnode, draw=grayDark!70, fill=white, font=\scriptsize,
                  text width=34mm, minimum height=8mm},
    tmine/.style={dnode, draw=redDark, line width=1pt, fill=flow,
                  font=\scriptsize, text width=34mm, minimum height=8mm},
    %
    % --- grouping panels ---
    panelA/.style={rounded corners=9pt, draw=blue!45,       dashed,
                   line width=1pt, fill=panelBlue},
    panelB/.style={rounded corners=9pt, draw=green!45!black, dashed,
                   line width=1pt, fill=panelGreen},
    panelC/.style={rounded corners=9pt, draw=redDark!55,     dashed,
                   line width=1pt, fill=panelWarm},
    %
    % --- annotation ---
    note/.style={font=\scriptsize, align=center, text=grayDark},
    tag/.style={font=\scriptsize, align=center, fill=white, inner sep=1.5pt},
    %
    % --- legend ---
    chip/.style={rounded corners=3pt, minimum width=.65cm,
                 minimum height=.3cm, inner sep=0pt},
    chiptext/.style={anchor=west, font=\scriptsize},
    legendbox/.style={rounded corners=7pt, draw=black!12, fill=white},
}

% Numbered section badge, used to title a panel.
% #1 stroke/fill colour, #2 node name, #3 the number, #4 anchor node,
% #5 the title text.
\newcommand{\panelbadge}[5]{%
  \node[circle, draw=#1, fill=#1, text=white, minimum size=.62cm,
        font=\small\bfseries] (#2) at ($(#4.north west)+(0.72,-0.58)$) {#3};
  \node[anchor=west, font=\bfseries, text=#1, right=2mm of #2] {#5};
}

% A legend swatch chained to the right of the previous one.
% #1 draw colour, #2 fill colour, #3 node name, #4 previous text node,
% #5 label
\newcommand{\legendchip}[5]{%
  \node[chip, draw=#1, fill=#2, right=8mm of #4] (#3) {};
  \node[chiptext, right=2mm of #3] (#3t) {#5};
}

%   \begin{document} \begin{tikzpicture} ... \end{tikzpicture} \end{document}
%
% The palette is the one introduced by diagram-orbis.tex. That file
% still carries its own inline copy of these definitions; if you ever
% touch the palette, change it in both places or switch orbis over to
% % ============================================================
% Shared preamble for every standalone diagram in this folder.
%
% Usage, from a file in this same directory:
%
%   \documentclass[tikz,border=8pt]{standalone}
%   % ============================================================
% Shared preamble for every standalone diagram in this folder.
%
% Usage, from a file in this same directory:
%
%   \documentclass[tikz,border=8pt]{standalone}
%   \input{diagramstyle}
%   \begin{document} \begin{tikzpicture} ... \end{tikzpicture} \end{document}
%
% The palette is the one introduced by diagram-orbis.tex. That file
% still carries its own inline copy of these definitions; if you ever
% touch the palette, change it in both places or switch orbis over to
% \input{diagramstyle} as well.
% ============================================================

\usepackage{amsmath,amssymb}
\usepackage{tikz}
\usetikzlibrary{arrows.meta, positioning, calc, fit, backgrounds,
                shapes.geometric, matrix, decorations.pathreplacing}

% ---------- fills -------------------------------------------
% Panel fills are OPAQUE pre-blended colours rather than transparent
% ones, so a line can be given a halo in exactly the panel colour.
% That is what turns an unavoidable crossing into a clean "hop".
\definecolor{panelBlue}{RGB}{248,251,255}
\definecolor{panelGreen}{RGB}{249,253,249}
\definecolor{panelWarm}{RGB}{255,251,247}

\definecolor{semantic}{RGB}{229,224,250}
\definecolor{detail}{RGB}{218,246,228}
\definecolor{latent}{RGB}{255,244,210}
\definecolor{vq}{RGB}{244,232,255}
\definecolor{flow}{RGB}{255,231,222}
\definecolor{attention}{RGB}{224,239,255}
\definecolor{decoder}{RGB}{231,236,255}
\definecolor{neutral}{RGB}{243,244,246}

% ---------- strokes -----------------------------------------
\definecolor{blueDark}{RGB}{25,73,115}
\definecolor{greenDark}{RGB}{35,120,78}
\definecolor{purpleDark}{RGB}{105,70,145}
\definecolor{redDark}{RGB}{145,65,55}
\definecolor{grayDark}{RGB}{78,82,90}
\definecolor{orangeDark}{RGB}{165,95,40}

% ---------- styles ------------------------------------------
\tikzset{
    >=Latex,
    % --- connectors ---
    arrow/.style={->, line width=0.8pt, shorten >=1.5pt, shorten <=1.5pt},
    thinarrow/.style={->, line width=0.6pt, shorten >=1.2pt, shorten <=1.2pt},
    plainline/.style={line width=0.8pt},
    dashedarrow/.style={->, dashed, line width=0.8pt,
                        shorten >=1.5pt, shorten <=1.5pt},
    xlink/.style={->, dashed, line width=0.9pt, redDark!70,
                  shorten >=1.5pt, shorten <=1.5pt},
    hop/.style={preaction={draw, line width=3.4pt, white}},
    junction/.style={circle, fill=black, inner sep=0pt, minimum size=3.4pt},
    %
    % --- boxes: every one carries generous inner sep so text never
    %     touches its own border ---
    dnode/.style={rounded corners=5pt, align=center, font=\small,
                  inner xsep=5pt, inner ysep=4pt, minimum height=9mm},
    io/.style     ={dnode, draw=grayDark,      fill=neutral},
    proc/.style   ={dnode, draw=blue!60!black, fill=attention},
    latentn/.style={dnode, draw=orange!65!black, fill=latent},
    learn/.style  ={dnode, draw=redDark,       fill=flow},
    attn/.style   ={dnode, draw=blue!60!black, fill=attention},
    head/.style   ={dnode, draw=greenDark,     fill=detail},
    cond/.style   ={dnode, draw=purpleDark,    fill=vq},
    seqn/.style   ={dnode, draw=blue!55!black, fill=decoder},
    frozen/.style ={dnode, draw=blue!55!black, fill=decoder},
    %
    % --- tree nodes ---
    troot/.style={dnode, draw=grayDark, fill=neutral, font=\small\bfseries,
                  text width=34mm},
    tcat/.style ={dnode, draw=blue!60!black, fill=attention,
                  font=\small\bfseries, text width=30mm},
    tleaf/.style={dnode, draw=grayDark!70, fill=white, font=\scriptsize,
                  text width=34mm, minimum height=8mm},
    tmine/.style={dnode, draw=redDark, line width=1pt, fill=flow,
                  font=\scriptsize, text width=34mm, minimum height=8mm},
    %
    % --- grouping panels ---
    panelA/.style={rounded corners=9pt, draw=blue!45,       dashed,
                   line width=1pt, fill=panelBlue},
    panelB/.style={rounded corners=9pt, draw=green!45!black, dashed,
                   line width=1pt, fill=panelGreen},
    panelC/.style={rounded corners=9pt, draw=redDark!55,     dashed,
                   line width=1pt, fill=panelWarm},
    %
    % --- annotation ---
    note/.style={font=\scriptsize, align=center, text=grayDark},
    tag/.style={font=\scriptsize, align=center, fill=white, inner sep=1.5pt},
    %
    % --- legend ---
    chip/.style={rounded corners=3pt, minimum width=.65cm,
                 minimum height=.3cm, inner sep=0pt},
    chiptext/.style={anchor=west, font=\scriptsize},
    legendbox/.style={rounded corners=7pt, draw=black!12, fill=white},
}

% Numbered section badge, used to title a panel.
% #1 stroke/fill colour, #2 node name, #3 the number, #4 anchor node,
% #5 the title text.
\newcommand{\panelbadge}[5]{%
  \node[circle, draw=#1, fill=#1, text=white, minimum size=.62cm,
        font=\small\bfseries] (#2) at ($(#4.north west)+(0.72,-0.58)$) {#3};
  \node[anchor=west, font=\bfseries, text=#1, right=2mm of #2] {#5};
}

% A legend swatch chained to the right of the previous one.
% #1 draw colour, #2 fill colour, #3 node name, #4 previous text node,
% #5 label
\newcommand{\legendchip}[5]{%
  \node[chip, draw=#1, fill=#2, right=8mm of #4] (#3) {};
  \node[chiptext, right=2mm of #3] (#3t) {#5};
}

%   \begin{document} \begin{tikzpicture} ... \end{tikzpicture} \end{document}
%
% The palette is the one introduced by diagram-orbis.tex. That file
% still carries its own inline copy of these definitions; if you ever
% touch the palette, change it in both places or switch orbis over to
% % ============================================================
% Shared preamble for every standalone diagram in this folder.
%
% Usage, from a file in this same directory:
%
%   \documentclass[tikz,border=8pt]{standalone}
%   \input{diagramstyle}
%   \begin{document} \begin{tikzpicture} ... \end{tikzpicture} \end{document}
%
% The palette is the one introduced by diagram-orbis.tex. That file
% still carries its own inline copy of these definitions; if you ever
% touch the palette, change it in both places or switch orbis over to
% \input{diagramstyle} as well.
% ============================================================

\usepackage{amsmath,amssymb}
\usepackage{tikz}
\usetikzlibrary{arrows.meta, positioning, calc, fit, backgrounds,
                shapes.geometric, matrix, decorations.pathreplacing}

% ---------- fills -------------------------------------------
% Panel fills are OPAQUE pre-blended colours rather than transparent
% ones, so a line can be given a halo in exactly the panel colour.
% That is what turns an unavoidable crossing into a clean "hop".
\definecolor{panelBlue}{RGB}{248,251,255}
\definecolor{panelGreen}{RGB}{249,253,249}
\definecolor{panelWarm}{RGB}{255,251,247}

\definecolor{semantic}{RGB}{229,224,250}
\definecolor{detail}{RGB}{218,246,228}
\definecolor{latent}{RGB}{255,244,210}
\definecolor{vq}{RGB}{244,232,255}
\definecolor{flow}{RGB}{255,231,222}
\definecolor{attention}{RGB}{224,239,255}
\definecolor{decoder}{RGB}{231,236,255}
\definecolor{neutral}{RGB}{243,244,246}

% ---------- strokes -----------------------------------------
\definecolor{blueDark}{RGB}{25,73,115}
\definecolor{greenDark}{RGB}{35,120,78}
\definecolor{purpleDark}{RGB}{105,70,145}
\definecolor{redDark}{RGB}{145,65,55}
\definecolor{grayDark}{RGB}{78,82,90}
\definecolor{orangeDark}{RGB}{165,95,40}

% ---------- styles ------------------------------------------
\tikzset{
    >=Latex,
    % --- connectors ---
    arrow/.style={->, line width=0.8pt, shorten >=1.5pt, shorten <=1.5pt},
    thinarrow/.style={->, line width=0.6pt, shorten >=1.2pt, shorten <=1.2pt},
    plainline/.style={line width=0.8pt},
    dashedarrow/.style={->, dashed, line width=0.8pt,
                        shorten >=1.5pt, shorten <=1.5pt},
    xlink/.style={->, dashed, line width=0.9pt, redDark!70,
                  shorten >=1.5pt, shorten <=1.5pt},
    hop/.style={preaction={draw, line width=3.4pt, white}},
    junction/.style={circle, fill=black, inner sep=0pt, minimum size=3.4pt},
    %
    % --- boxes: every one carries generous inner sep so text never
    %     touches its own border ---
    dnode/.style={rounded corners=5pt, align=center, font=\small,
                  inner xsep=5pt, inner ysep=4pt, minimum height=9mm},
    io/.style     ={dnode, draw=grayDark,      fill=neutral},
    proc/.style   ={dnode, draw=blue!60!black, fill=attention},
    latentn/.style={dnode, draw=orange!65!black, fill=latent},
    learn/.style  ={dnode, draw=redDark,       fill=flow},
    attn/.style   ={dnode, draw=blue!60!black, fill=attention},
    head/.style   ={dnode, draw=greenDark,     fill=detail},
    cond/.style   ={dnode, draw=purpleDark,    fill=vq},
    seqn/.style   ={dnode, draw=blue!55!black, fill=decoder},
    frozen/.style ={dnode, draw=blue!55!black, fill=decoder},
    %
    % --- tree nodes ---
    troot/.style={dnode, draw=grayDark, fill=neutral, font=\small\bfseries,
                  text width=34mm},
    tcat/.style ={dnode, draw=blue!60!black, fill=attention,
                  font=\small\bfseries, text width=30mm},
    tleaf/.style={dnode, draw=grayDark!70, fill=white, font=\scriptsize,
                  text width=34mm, minimum height=8mm},
    tmine/.style={dnode, draw=redDark, line width=1pt, fill=flow,
                  font=\scriptsize, text width=34mm, minimum height=8mm},
    %
    % --- grouping panels ---
    panelA/.style={rounded corners=9pt, draw=blue!45,       dashed,
                   line width=1pt, fill=panelBlue},
    panelB/.style={rounded corners=9pt, draw=green!45!black, dashed,
                   line width=1pt, fill=panelGreen},
    panelC/.style={rounded corners=9pt, draw=redDark!55,     dashed,
                   line width=1pt, fill=panelWarm},
    %
    % --- annotation ---
    note/.style={font=\scriptsize, align=center, text=grayDark},
    tag/.style={font=\scriptsize, align=center, fill=white, inner sep=1.5pt},
    %
    % --- legend ---
    chip/.style={rounded corners=3pt, minimum width=.65cm,
                 minimum height=.3cm, inner sep=0pt},
    chiptext/.style={anchor=west, font=\scriptsize},
    legendbox/.style={rounded corners=7pt, draw=black!12, fill=white},
}

% Numbered section badge, used to title a panel.
% #1 stroke/fill colour, #2 node name, #3 the number, #4 anchor node,
% #5 the title text.
\newcommand{\panelbadge}[5]{%
  \node[circle, draw=#1, fill=#1, text=white, minimum size=.62cm,
        font=\small\bfseries] (#2) at ($(#4.north west)+(0.72,-0.58)$) {#3};
  \node[anchor=west, font=\bfseries, text=#1, right=2mm of #2] {#5};
}

% A legend swatch chained to the right of the previous one.
% #1 draw colour, #2 fill colour, #3 node name, #4 previous text node,
% #5 label
\newcommand{\legendchip}[5]{%
  \node[chip, draw=#1, fill=#2, right=8mm of #4] (#3) {};
  \node[chiptext, right=2mm of #3] (#3t) {#5};
}
 as well.
% ============================================================

\usepackage{amsmath,amssymb}
\usepackage{tikz}
\usetikzlibrary{arrows.meta, positioning, calc, fit, backgrounds,
                shapes.geometric, matrix, decorations.pathreplacing}

% ---------- fills -------------------------------------------
% Panel fills are OPAQUE pre-blended colours rather than transparent
% ones, so a line can be given a halo in exactly the panel colour.
% That is what turns an unavoidable crossing into a clean "hop".
\definecolor{panelBlue}{RGB}{248,251,255}
\definecolor{panelGreen}{RGB}{249,253,249}
\definecolor{panelWarm}{RGB}{255,251,247}

\definecolor{semantic}{RGB}{229,224,250}
\definecolor{detail}{RGB}{218,246,228}
\definecolor{latent}{RGB}{255,244,210}
\definecolor{vq}{RGB}{244,232,255}
\definecolor{flow}{RGB}{255,231,222}
\definecolor{attention}{RGB}{224,239,255}
\definecolor{decoder}{RGB}{231,236,255}
\definecolor{neutral}{RGB}{243,244,246}

% ---------- strokes -----------------------------------------
\definecolor{blueDark}{RGB}{25,73,115}
\definecolor{greenDark}{RGB}{35,120,78}
\definecolor{purpleDark}{RGB}{105,70,145}
\definecolor{redDark}{RGB}{145,65,55}
\definecolor{grayDark}{RGB}{78,82,90}
\definecolor{orangeDark}{RGB}{165,95,40}

% ---------- styles ------------------------------------------
\tikzset{
    >=Latex,
    % --- connectors ---
    arrow/.style={->, line width=0.8pt, shorten >=1.5pt, shorten <=1.5pt},
    thinarrow/.style={->, line width=0.6pt, shorten >=1.2pt, shorten <=1.2pt},
    plainline/.style={line width=0.8pt},
    dashedarrow/.style={->, dashed, line width=0.8pt,
                        shorten >=1.5pt, shorten <=1.5pt},
    xlink/.style={->, dashed, line width=0.9pt, redDark!70,
                  shorten >=1.5pt, shorten <=1.5pt},
    hop/.style={preaction={draw, line width=3.4pt, white}},
    junction/.style={circle, fill=black, inner sep=0pt, minimum size=3.4pt},
    %
    % --- boxes: every one carries generous inner sep so text never
    %     touches its own border ---
    dnode/.style={rounded corners=5pt, align=center, font=\small,
                  inner xsep=5pt, inner ysep=4pt, minimum height=9mm},
    io/.style     ={dnode, draw=grayDark,      fill=neutral},
    proc/.style   ={dnode, draw=blue!60!black, fill=attention},
    latentn/.style={dnode, draw=orange!65!black, fill=latent},
    learn/.style  ={dnode, draw=redDark,       fill=flow},
    attn/.style   ={dnode, draw=blue!60!black, fill=attention},
    head/.style   ={dnode, draw=greenDark,     fill=detail},
    cond/.style   ={dnode, draw=purpleDark,    fill=vq},
    seqn/.style   ={dnode, draw=blue!55!black, fill=decoder},
    frozen/.style ={dnode, draw=blue!55!black, fill=decoder},
    %
    % --- tree nodes ---
    troot/.style={dnode, draw=grayDark, fill=neutral, font=\small\bfseries,
                  text width=34mm},
    tcat/.style ={dnode, draw=blue!60!black, fill=attention,
                  font=\small\bfseries, text width=30mm},
    tleaf/.style={dnode, draw=grayDark!70, fill=white, font=\scriptsize,
                  text width=34mm, minimum height=8mm},
    tmine/.style={dnode, draw=redDark, line width=1pt, fill=flow,
                  font=\scriptsize, text width=34mm, minimum height=8mm},
    %
    % --- grouping panels ---
    panelA/.style={rounded corners=9pt, draw=blue!45,       dashed,
                   line width=1pt, fill=panelBlue},
    panelB/.style={rounded corners=9pt, draw=green!45!black, dashed,
                   line width=1pt, fill=panelGreen},
    panelC/.style={rounded corners=9pt, draw=redDark!55,     dashed,
                   line width=1pt, fill=panelWarm},
    %
    % --- annotation ---
    note/.style={font=\scriptsize, align=center, text=grayDark},
    tag/.style={font=\scriptsize, align=center, fill=white, inner sep=1.5pt},
    %
    % --- legend ---
    chip/.style={rounded corners=3pt, minimum width=.65cm,
                 minimum height=.3cm, inner sep=0pt},
    chiptext/.style={anchor=west, font=\scriptsize},
    legendbox/.style={rounded corners=7pt, draw=black!12, fill=white},
}

% Numbered section badge, used to title a panel.
% #1 stroke/fill colour, #2 node name, #3 the number, #4 anchor node,
% #5 the title text.
\newcommand{\panelbadge}[5]{%
  \node[circle, draw=#1, fill=#1, text=white, minimum size=.62cm,
        font=\small\bfseries] (#2) at ($(#4.north west)+(0.72,-0.58)$) {#3};
  \node[anchor=west, font=\bfseries, text=#1, right=2mm of #2] {#5};
}

% A legend swatch chained to the right of the previous one.
% #1 draw colour, #2 fill colour, #3 node name, #4 previous text node,
% #5 label
\newcommand{\legendchip}[5]{%
  \node[chip, draw=#1, fill=#2, right=8mm of #4] (#3) {};
  \node[chiptext, right=2mm of #3] (#3t) {#5};
}
 as well.
% ============================================================

\usepackage{amsmath,amssymb}
\usepackage{tikz}
\usetikzlibrary{arrows.meta, positioning, calc, fit, backgrounds,
                shapes.geometric, matrix, decorations.pathreplacing}

% ---------- fills -------------------------------------------
% Panel fills are OPAQUE pre-blended colours rather than transparent
% ones, so a line can be given a halo in exactly the panel colour.
% That is what turns an unavoidable crossing into a clean "hop".
\definecolor{panelBlue}{RGB}{248,251,255}
\definecolor{panelGreen}{RGB}{249,253,249}
\definecolor{panelWarm}{RGB}{255,251,247}

\definecolor{semantic}{RGB}{229,224,250}
\definecolor{detail}{RGB}{218,246,228}
\definecolor{latent}{RGB}{255,244,210}
\definecolor{vq}{RGB}{244,232,255}
\definecolor{flow}{RGB}{255,231,222}
\definecolor{attention}{RGB}{224,239,255}
\definecolor{decoder}{RGB}{231,236,255}
\definecolor{neutral}{RGB}{243,244,246}

% ---------- strokes -----------------------------------------
\definecolor{blueDark}{RGB}{25,73,115}
\definecolor{greenDark}{RGB}{35,120,78}
\definecolor{purpleDark}{RGB}{105,70,145}
\definecolor{redDark}{RGB}{145,65,55}
\definecolor{grayDark}{RGB}{78,82,90}
\definecolor{orangeDark}{RGB}{165,95,40}

% ---------- styles ------------------------------------------
\tikzset{
    >=Latex,
    % --- connectors ---
    arrow/.style={->, line width=0.8pt, shorten >=1.5pt, shorten <=1.5pt},
    thinarrow/.style={->, line width=0.6pt, shorten >=1.2pt, shorten <=1.2pt},
    plainline/.style={line width=0.8pt},
    dashedarrow/.style={->, dashed, line width=0.8pt,
                        shorten >=1.5pt, shorten <=1.5pt},
    xlink/.style={->, dashed, line width=0.9pt, redDark!70,
                  shorten >=1.5pt, shorten <=1.5pt},
    hop/.style={preaction={draw, line width=3.4pt, white}},
    junction/.style={circle, fill=black, inner sep=0pt, minimum size=3.4pt},
    %
    % --- boxes: every one carries generous inner sep so text never
    %     touches its own border ---
    dnode/.style={rounded corners=5pt, align=center, font=\small,
                  inner xsep=5pt, inner ysep=4pt, minimum height=9mm},
    io/.style     ={dnode, draw=grayDark,      fill=neutral},
    proc/.style   ={dnode, draw=blue!60!black, fill=attention},
    latentn/.style={dnode, draw=orange!65!black, fill=latent},
    learn/.style  ={dnode, draw=redDark,       fill=flow},
    attn/.style   ={dnode, draw=blue!60!black, fill=attention},
    head/.style   ={dnode, draw=greenDark,     fill=detail},
    cond/.style   ={dnode, draw=purpleDark,    fill=vq},
    seqn/.style   ={dnode, draw=blue!55!black, fill=decoder},
    frozen/.style ={dnode, draw=blue!55!black, fill=decoder},
    %
    % --- tree nodes ---
    troot/.style={dnode, draw=grayDark, fill=neutral, font=\small\bfseries,
                  text width=34mm},
    tcat/.style ={dnode, draw=blue!60!black, fill=attention,
                  font=\small\bfseries, text width=30mm},
    tleaf/.style={dnode, draw=grayDark!70, fill=white, font=\scriptsize,
                  text width=34mm, minimum height=8mm},
    tmine/.style={dnode, draw=redDark, line width=1pt, fill=flow,
                  font=\scriptsize, text width=34mm, minimum height=8mm},
    %
    % --- grouping panels ---
    panelA/.style={rounded corners=9pt, draw=blue!45,       dashed,
                   line width=1pt, fill=panelBlue},
    panelB/.style={rounded corners=9pt, draw=green!45!black, dashed,
                   line width=1pt, fill=panelGreen},
    panelC/.style={rounded corners=9pt, draw=redDark!55,     dashed,
                   line width=1pt, fill=panelWarm},
    %
    % --- annotation ---
    note/.style={font=\scriptsize, align=center, text=grayDark},
    tag/.style={font=\scriptsize, align=center, fill=white, inner sep=1.5pt},
    %
    % --- legend ---
    chip/.style={rounded corners=3pt, minimum width=.65cm,
                 minimum height=.3cm, inner sep=0pt},
    chiptext/.style={anchor=west, font=\scriptsize},
    legendbox/.style={rounded corners=7pt, draw=black!12, fill=white},
}

% Numbered section badge, used to title a panel.
% #1 stroke/fill colour, #2 node name, #3 the number, #4 anchor node,
% #5 the title text.
\newcommand{\panelbadge}[5]{%
  \node[circle, draw=#1, fill=#1, text=white, minimum size=.62cm,
        font=\small\bfseries] (#2) at ($(#4.north west)+(0.72,-0.58)$) {#3};
  \node[anchor=west, font=\bfseries, text=#1, right=2mm of #2] {#5};
}

% A legend swatch chained to the right of the previous one.
% #1 draw colour, #2 fill colour, #3 node name, #4 previous text node,
% #5 label
\newcommand{\legendchip}[5]{%
  \node[chip, draw=#1, fill=#2, right=8mm of #4] (#3) {};
  \node[chiptext, right=2mm of #3] (#3t) {#5};
}
 as well.
% ============================================================

\usepackage{amsmath,amssymb}
\usepackage{tikz}
\usetikzlibrary{arrows.meta, positioning, calc, fit, backgrounds,
                shapes.geometric, matrix, decorations.pathreplacing}

% ---------- fills -------------------------------------------
% Panel fills are OPAQUE pre-blended colours rather than transparent
% ones, so a line can be given a halo in exactly the panel colour.
% That is what turns an unavoidable crossing into a clean "hop".
\definecolor{panelBlue}{RGB}{248,251,255}
\definecolor{panelGreen}{RGB}{249,253,249}
\definecolor{panelWarm}{RGB}{255,251,247}

\definecolor{semantic}{RGB}{229,224,250}
\definecolor{detail}{RGB}{218,246,228}
\definecolor{latent}{RGB}{255,244,210}
\definecolor{vq}{RGB}{244,232,255}
\definecolor{flow}{RGB}{255,231,222}
\definecolor{attention}{RGB}{224,239,255}
\definecolor{decoder}{RGB}{231,236,255}
\definecolor{neutral}{RGB}{243,244,246}

% ---------- strokes -----------------------------------------
\definecolor{blueDark}{RGB}{25,73,115}
\definecolor{greenDark}{RGB}{35,120,78}
\definecolor{purpleDark}{RGB}{105,70,145}
\definecolor{redDark}{RGB}{145,65,55}
\definecolor{grayDark}{RGB}{78,82,90}
\definecolor{orangeDark}{RGB}{165,95,40}

% ---------- styles ------------------------------------------
\tikzset{
    >=Latex,
    % --- connectors ---
    arrow/.style={->, line width=0.8pt, shorten >=1.5pt, shorten <=1.5pt},
    thinarrow/.style={->, line width=0.6pt, shorten >=1.2pt, shorten <=1.2pt},
    plainline/.style={line width=0.8pt},
    dashedarrow/.style={->, dashed, line width=0.8pt,
                        shorten >=1.5pt, shorten <=1.5pt},
    xlink/.style={->, dashed, line width=0.9pt, redDark!70,
                  shorten >=1.5pt, shorten <=1.5pt},
    hop/.style={preaction={draw, line width=3.4pt, white}},
    junction/.style={circle, fill=black, inner sep=0pt, minimum size=3.4pt},
    %
    % --- boxes: every one carries generous inner sep so text never
    %     touches its own border ---
    dnode/.style={rounded corners=5pt, align=center, font=\small,
                  inner xsep=5pt, inner ysep=4pt, minimum height=9mm},
    io/.style     ={dnode, draw=grayDark,      fill=neutral},
    proc/.style   ={dnode, draw=blue!60!black, fill=attention},
    latentn/.style={dnode, draw=orange!65!black, fill=latent},
    learn/.style  ={dnode, draw=redDark,       fill=flow},
    attn/.style   ={dnode, draw=blue!60!black, fill=attention},
    head/.style   ={dnode, draw=greenDark,     fill=detail},
    cond/.style   ={dnode, draw=purpleDark,    fill=vq},
    seqn/.style   ={dnode, draw=blue!55!black, fill=decoder},
    frozen/.style ={dnode, draw=blue!55!black, fill=decoder},
    %
    % --- tree nodes ---
    troot/.style={dnode, draw=grayDark, fill=neutral, font=\small\bfseries,
                  text width=34mm},
    tcat/.style ={dnode, draw=blue!60!black, fill=attention,
                  font=\small\bfseries, text width=30mm},
    tleaf/.style={dnode, draw=grayDark!70, fill=white, font=\scriptsize,
                  text width=34mm, minimum height=8mm},
    tmine/.style={dnode, draw=redDark, line width=1pt, fill=flow,
                  font=\scriptsize, text width=34mm, minimum height=8mm},
    %
    % --- grouping panels ---
    panelA/.style={rounded corners=9pt, draw=blue!45,       dashed,
                   line width=1pt, fill=panelBlue},
    panelB/.style={rounded corners=9pt, draw=green!45!black, dashed,
                   line width=1pt, fill=panelGreen},
    panelC/.style={rounded corners=9pt, draw=redDark!55,     dashed,
                   line width=1pt, fill=panelWarm},
    %
    % --- annotation ---
    note/.style={font=\scriptsize, align=center, text=grayDark},
    tag/.style={font=\scriptsize, align=center, fill=white, inner sep=1.5pt},
    %
    % --- legend ---
    chip/.style={rounded corners=3pt, minimum width=.65cm,
                 minimum height=.3cm, inner sep=0pt},
    chiptext/.style={anchor=west, font=\scriptsize},
    legendbox/.style={rounded corners=7pt, draw=black!12, fill=white},
}

% Numbered section badge, used to title a panel.
% #1 stroke/fill colour, #2 node name, #3 the number, #4 anchor node,
% #5 the title text.
\newcommand{\panelbadge}[5]{%
  \node[circle, draw=#1, fill=#1, text=white, minimum size=.62cm,
        font=\small\bfseries] (#2) at ($(#4.north west)+(0.72,-0.58)$) {#3};
  \node[anchor=west, font=\bfseries, text=#1, right=2mm of #2] {#5};
}

% A legend swatch chained to the right of the previous one.
% #1 draw colour, #2 fill colour, #3 node name, #4 previous text node,
% #5 label
\newcommand{\legendchip}[5]{%
  \node[chip, draw=#1, fill=#2, right=8mm of #4] (#3) {};
  \node[chiptext, right=2mm of #3] (#3t) {#5};
}
